\section{Conclusion and Limitations}
\label{sec:discussion}

ToxicBench measures how data agents respond when a successful tool call returns wrong evidence. Across frameworks and models, poisoning reduces task success and can lead agents to repeat a corrupted value or label. The verification experiments add two findings. An ordinary second pass can recover many numerical tasks, while generic guard prompts do not improve on that control. When later observations are also poisoned, agents may check the evidence and still adopt the wrong answer. Together, the task pairs, trajectory metrics, and verification controls make these failures visible and provide a common setting for testing better evidence use.

The current benchmark focuses on tabular tasks and controlled fault injection. Verification routes share source tables and backends, and adapters expose different tool boundaries. The scorer also has measurement limits: the development audit gives low BCR recall, and TSR, PAR, and VPA have not received separate held-out human validation. Task bootstrap intervals do not include scorer error. Future work should add an independent answer-level audit and extend the tasks to live services, larger databases, and checks that use separate data sources.
